\section{Experiments}\label{sec:exp}

This section measures the implementation of \cref{sec:impl}.
The goals are modest and specific: to check that the kernel encoding costs the same as coefficient-form folding at identical engineering (as predicted by \cref{rem:basefold}), to locate where the time goes, and to quantify the two protocol parameters $\ell$ and $\rho$.
We do not claim competitive absolute performance: the code is single-threaded and uses no SIMD, no Merkle multiproofs and no hash-specific batching.

\subsection{Setup}\label{sec:exp:setup}

\paragraph{Machine.}
All runs use one core of an Intel Xeon processor at $2.10$\,GHz (cloud virtual machine, $2$ vCPUs, $7$\,GB RAM), Rust $1.95.0$, release profile with fat LTO and one codegen unit, and no \texttt{target-cpu} flag.
Times are medians over $11$ runs for $n \le 18$, $5$ for $n = 20$ and $3$ for $n = 22$ and for \cref{tab:stop,tab:rate}; verification times are medians over $21$ runs.
The machine is a shared virtual machine, and timings of identical work vary by up to about $20\%$ between runs (see \cref{sec:exp:scaling}).

\paragraph{Parameters.}
Unless stated otherwise: rate $\rho = 1/4$ ($R = 2$), number of folding rounds $\ell = n - 4$ (final table of $16$ elements of $\K$), and $\kappa = 148$ queries.
By \cref{thm:sound} with $\delta = (1-\rho)/2 = 3/8$, the query term is $(5/8)^{148} < 2^{-100}$, and $\varepsilon_{\mathrm{fold}} + 2\ell/|\K| < 2^{-100}$ for all sizes tested ($M \le 2^{24}$, $|\K| \approx 2^{128}$).
Tables $f$ are uniformly random over $\mathbb{F}_p$ and the point $z$ is uniformly random in $\K^n$.

\paragraph{Baseline.}
To isolate the effect of the encoding, the library has a second mode that differs from $\Pieval$ in exactly two places: the table is converted to the monomial coefficients of $\tilde f$ by M\"obius inversion (\cref{lem:mobius}) instead of the transform of \cref{cor:transforms}, and words are folded with the classical FRI fold $\mathrm{cfold}_\theta(w) = w_e + \theta\,w_o$ (\cref{def:cfold}) instead of the kernel fold (\cref{rem:coeff-form}).
This is the coefficient-form folding used in Reed--Solomon instantiations of BaseFold and in DeepFold; everything else (field, NTT, Merkle trees, transcript, sumcheck, query phase) is shared.
The baseline passes the same completeness and tampering tests.
It is \emph{not} the BaseFold or DeepFold code base, and \cref{tab:scaling} should be read as a comparison of encodings, not of implementations.

\subsection{Kernel versus coefficient-form encoding}\label{sec:exp:scaling}

\begin{table}[t]
\centering
\small
\begin{tabular}{r rr rr rr r}
\toprule
 & \multicolumn{2}{c}{Commit (ms)} & \multicolumn{2}{c}{Open (ms)} & \multicolumn{2}{c}{Verify (ms)} & Proof \\
$n$ & kernel & coeff. & kernel & coeff. & kernel & coeff. & (KiB) \\
\midrule
12 & 2.5 & 2.3 & 2.6 & 2.6 & 2.3 & 2.3 & 387 \\
14 & 10.2 & 9.8 & 10.4 & 10.3 & 3.1 & 3.1 & 531 \\
16 & 44.7 & 36.5 & 42.1 & 34.8 & 3.1 & 3.5 & 693 \\
18 & 158.8 & 176.4 & 173.6 & 171.7 & 5.0 & 4.1 & 873 \\
20 & 821.5 & 779.0 & 734.8 & 751.8 & 6.1 & 6.1 & 1072 \\
22 & 4139.7 & 3998.1 & 3354.9 & 3401.3 & 6.1 & 7.2 & 1290 \\
\bottomrule
\end{tabular}
\caption{Kernel encoding ($\Pieval$) versus the coefficient-form baseline, single-threaded, $\rho = 1/4$, $\ell = n-4$, $\kappa = 148$. Proof sizes are identical for both encodings.}
\label{tab:scaling}
\end{table}

\Cref{tab:scaling} reports commit, open and verify times for $n = 12,\dots,22$ ($N = 2^{12},\dots,2^{22}$).
Proof sizes are identical because the proofs have the same shape.
The timing differences between the two encodings are of the same size as the differences in steps whose code is identical in both modes: in \cref{tab:breakdown}, the NTT and Merkle columns differ by up to about $20\%$ between modes although they perform exactly the same computation.
Moreover, in an earlier full run of the same benchmark, the sign of the difference in commit time reversed at $n = 18$ and $n = 20$.
We conclude that the two encodings are indistinguishable at the resolution of this machine.
This is the outcome predicted by \cref{rem:basefold}: the kernel transform and M\"obius inversion both cost $\tfrac{n}{2}N$ additions, and in line form (\cref{lem:line}) the kernel fold costs one multiplication by $r$ per output element, the same as the classical fold.
The kernel encoding therefore buys its structural property (\cref{thm:fold-dictionary}) at no measurable cost.

\subsection{Where the time goes}\label{sec:exp:breakdown}

\begin{table}[t]
\centering
\small
\begin{tabular}{r rr rr rr}
\toprule
 & \multicolumn{2}{c}{Transform (ms)} & \multicolumn{2}{c}{NTT (ms)} & \multicolumn{2}{c}{Merkle tree (ms)} \\
$n$ & kernel & coeff. & kernel & coeff. & kernel & coeff. \\
\midrule
16 & 0.48 & 0.32 & 9.4 & 7.8 & 32.4 & 28.4 \\
18 & 2.68 & 2.61 & 40.6 & 46.3 & 104.6 & 128.6 \\
20 & 15.2 & 13.6 & 210.2 & 227.5 & 521.2 & 509.1 \\
22 & 60.8 & 66.1 & 1476.1 & 1469.5 & 2391.2 & 2271.7 \\
\bottomrule
\end{tabular}
\caption{Components of the commitment ($M = 4N$), measured separately in the same runs as \cref{tab:scaling}. The NTT and Merkle columns execute identical code in both modes.}
\label{tab:breakdown}
\end{table}

\Cref{tab:breakdown} splits the commitment into its three steps.
The kernel transform accounts for $1$--$2\%$ of the commitment at every size, and takes the same time as M\"obius inversion (it performs $\tfrac{n}{2}N$ additions, against $\tfrac{M}{2}\log_2 M$ multiplications for the NTT on the $4\times$ larger domain).
The commitment is dominated by hashing: BLAKE3 over $M/2$ fibre leaves and $M/2 - 1$ internal nodes takes $1.6$--$3.5$ times as long as the NTT.
Opening costs about as much as committing: the folded words $w_1,\dots,w_{\ell-1}$ have total length below $M$ but live in $\K$, and each needs its own Merkle tree; the sumcheck ($O(N)$ multiplications in $\K$) is a minor term.

Verification takes $2$--$7$\,ms and grows slowly with $n$, as expected from its $O(\kappa \ell \log M)$ hashing cost.
Proof sizes ($0.4$--$1.3$\,MiB) are dominated by the $\kappa \ell$ authentication paths: we open one Merkle leaf per query and level and do not merge paths across queries.
Both the unique-decoding query count and the absence of path merging are deliberate simplifications; \cref{rem:params} discusses the former.

\subsection{The number of folding rounds}\label{sec:exp:stop}

\begin{table}[t]
\centering
\small
\begin{tabular}{r r rrr}
\toprule
$\ell$ & final table $N_\ell$ & open (ms) & verify (ms) & proof (KiB) \\
\midrule
8 & 4096 & 769.7 & 9.1 & 747 \\
10 & 1024 & 748.8 & 5.8 & 824 \\
12 & 256 & 746.9 & 4.5 & 918 \\
14 & 64 & 726.8 & 5.6 & 1003 \\
16 & 16 & 666.2 & 6.0 & 1072 \\
18 & 4 & 744.9 & 6.4 & 1123 \\
20 & 1 & 678.9 & 5.2 & 1156 \\
\bottomrule
\end{tabular}
\caption{Effect of the number of folding rounds $\ell$ (\cref{rem:stop}), kernel encoding, $n = 20$, $\rho = 1/4$, $\kappa = 148$.}
\label{tab:stop}
\end{table}

\Cref{tab:stop} varies $\ell$ at $n = 20$.
Each folding round adds $\kappa$ authentication paths to the proof, while stopping earlier adds $16 N_\ell$ bytes for the final table; in our range the paths dominate, so proofs shrink as $\ell$ decreases, from $1156$\,KiB at $\ell = n$ to $747$\,KiB at $\ell = n-12$.
Verification time is smallest around $\ell = 12$: for smaller $\ell$ the verifier's $\kappa$ evaluations of $G$, each of cost $O(N_\ell)$, take over.
Opening time is essentially independent of $\ell$, since the first folds dominate it.

\subsection{The rate}\label{sec:exp:rate}

\begin{table}[t]
\centering
\small
\begin{tabular}{r r r rrrr}
\toprule
$\rho$ & $\delta$ & $\kappa$ & commit (ms) & open (ms) & verify (ms) & proof (KiB) \\
\midrule
$1/2$ & $1/4$ & 241 & 422.2 & 405.3 & 7.9 & 1625 \\
$1/4$ & $3/8$ & 148 & 704.5 & 765.5 & 6.1 & 1072 \\
$1/8$ & $7/16$ & 121 & 1880.5 & 1521.2 & 5.3 & 937 \\
\bottomrule
\end{tabular}
\caption{Effect of the rate at $100$ bits of query security, $\delta = (1-\rho)/2$, kernel encoding, $n = 20$, $\ell = 16$.}
\label{tab:rate}
\end{table}

\Cref{tab:rate} fixes the query term of \cref{thm:sound} at $2^{-100}$ and varies the rate.
Halving the rate doubles the domain and increases the prover's work by a factor of $1.7$--$2.7$; it also reduces the number of queries, which shortens proofs and speeds up verification.
Within the unique-decoding analysis of this paper, $\rho = 1/4$ is a reasonable default.
